%------------------------------------------------------------------------------%
\section{Exercícios Propostos}

\begin{enumerate}
  \item Uma companhia de turismo cobra $R\$ 1200,00$ por um pacote turístico
    individual, mas dá um bom desconto quando os turistas viajam em grupo. Se um
    grupo tem $x$ pessoas, com $x \leq 70$, cada integrante paga apenas
  \[C(x)=1200-12x.\]
  Sabendo que a receita da empresa em cada grupo é dada pelo produto do valor pago por pessoa pelo número de membros do grupo.

\begin{enumerate}
\item escreva a função $R(x)$ que fornece a receita obtida pela empresa de turismo com o grupo de $x$ pessoas
\item trace o gráfico da função para $x \in [10,70]$
\item indique qual o tamanho do grupo fornece a maior receita
\end{enumerate}
\item Calcule
  \[\frac{f(x)-f(a)}{x-a}\]
para as funções e os valores de $a$ fornecidos a seguir. Simplifique os resultados e suponha sempre que os denominadores são diferentes de zero.
\begin{itemize}
\item[a)] $f(x)=3x+4, \hspace{1cm} a=2$
\item[b)] $f(x)=x^2+6, \hspace{1cm} a=4$
\end{itemize}

\item O gráfico da função $f(x)= \sqrt{-x^2+2x+3}$ é mostrado a seguir. Determine:

\begin{center}
  \includegraphics[width=0.5\textwidth]{questao_aula1-3}
\end{center}

\begin{itemize}
\item[a)] Os valores de $f(0),f(0,5)$ e $f(2)$
\item[b)] O domínio de $f$
\item[c)]O conjunto imagem de $f$
\item[d)] Os zeros de $f$
\item[e)] Os intervalos de crescimento e decrescimento
\item[f)] Os pontos de máximo e mínimo absoluto
\end{itemize}

\item Encontre o domínio das seguintes funções: \\
a) $f(x)= \frac{1}{x}+ \frac{1}{x+5}$ \hspace{1cm} b) $f(x)= \sqrt{x+2}+ \sqrt{x-5}$ \hspace{1cm} c) $f(x)=\frac{\sqrt{2x-3}}{x^2-5x+4}$ \hspace{1cm} d) $h(x)= \frac{1}{\sqrt[4]{x^2-5x}}$

\item Verifique algebricamente se as funções a seguir são pares, ímpares ou não possuem simetria. \\
a) $f(x)= x^2-3$ \hspace{1cm} b) $f(x)=-2x$ \hspace{1cm}  c) $f(x)=\frac{1}{x^2+1}$\hspace{1cm} d) $f(x)=2x^5-x^3+x$   \hspace{1cm} e) $y=x \sqrt{x}$ \hspace{1cm} f) $f(x)=2x-1$.

\item  Se $f(x)=x+2$, $g(x)=x^2-4$ e $h(x)=-2x+5$
\begin{enumerate}
\item Verifique se cada uma dessas funções são par, ímpar ou nem par nem ímpar
\item Verifique se essas funções são injetoras
\item Esboce o gráfico dessas funções, encontre o conjunto imagem e os intervalos de crescimento e decrescimento
\item Encontre o valor de  $\frac{g(1)-f(0)-h(3)}{h(1)}$
\item Calcule $g(h(2)))$
\item A partir do gráfico de $y=2^x$, esboce o gráfico de $y= 2^{f(x)}$
\end{enumerate}


\item (\textbf{Cercando um pasto}) Um fazendeiro pretende usar $400 m$ de cerca para delimitar uma área retangular que servirá de pasto. Responda os itens a seguir , lembrando que a área de um terreno retangular com largura $l$ e profundidade $p$ é dada por $l.p$, e que o perímetro desse terreno retangular é igual a $2l+2p$.
\begin{itemize}
\item[a)] Relacione a profundidade com a largura do pasto, considerando o uso dos $400 m$ de cerca.
\item[b)] Escreva uma função que forneça a área cercada em relação a largura do pasto.
\item[c)] Calcule a área do pasto, supondo que sua largura é $75m$. Faça o mesmo para uma largura de $150m$.
\item[d)] Determine o domínio da função área que você obteve.
\item[e)] Esboce o gráfico da função área.
\item[f)] Indique os intervalos nos quais a função é crescente e em quais ela é decrescente.
\item[g)] Com base no gráfico, indique se é possível cercar uma área de $12000 m^2$.
\item[h)] Ainda com base no gráfico, determine a maior área que pode ser cercada e as dimensões do terreno nesse caso.
\end{itemize}

\item A partir do gráfico de $f(x)=x^2$, esboce o gráfico das seguintes funções: \\
a) $g(x)= x^2-3$ \hspace{1cm} b) $h(x)= 3x^2-8$ \hspace{1cm} c)$p(x)=-x^2+6$ \hspace{1cm} d) $q(x)= \frac{1}{2}(x-2)^2$  \\ e) $u(x)=(\frac{1}{4}x)^2$.

\item A partir do gráfico de $f(x)=\sqrt[3]{x}$, esboce o gráfico das seguintes funções: \\
a) $g(x)= \sqrt[3]{x+3}$ \hspace{1cm} b) $h(x)= \sqrt[3]{3x}$ \hspace{1cm} c)$i(x)=-\sqrt[3]{x}$ \hspace{1cm} d) $j(x)= 3\sqrt[3]{x}$  \hspace{1cm} e) $k(x)=1-\sqrt[3]{x-1}$ \hspace{1cm} f)$l(x) =\sqrt[3]{2x}-1$.

\item Dadas as funções $f$ e $g$ a seguir, defina $f(g(x))$, $g(f(x))$, $f(f(x))$ e $g(g(x))$. \\
a) $f(x)=4x$, \hspace{0.1cm} $g(x)= \frac{x^2}{4}$ \hspace{1cm} b )$f(x)=\sqrt{x}$, \hspace{0.1cm} $g(x)= \frac{x}{3}$.

\item Dadas as funções $f(x)=\frac{1}{x-4}$ e $g(x)=x^2$,
\begin{itemize}
\item[a)]Defina $f(g(x))$ e $g(f(x))$ e seus domínios.
\item[b)] Calcule $f(g(-3))$ e $g(f(7))$.

\end{itemize}
\item Na cidade de Salicilina, o número de casos de determinada doença viral varia ao longo do ano, sendo dado aproximadamente por $d(t)=1,9245t^4-53,485t^3+450,18t^2-1081,5t+837,37$, em que $t$ é o mês do ano. Por sua vez, o custo em reais do tratamento dessa doença é dado pela função $c(x)=600x+12000$, em que $x$ é o número de pessoas infectadas.
\begin{itemize}
\item[a)] Determine $c(d(t))$.
\item[b)] Indique o que essa função composta representa.
\item[c)] Calcule o custo aproximado de tratamento dos salicilinenses infectados no mês de abril. (Use uma calculadora).
\end{itemize}
\end{enumerate}

%------------------------------------------------------------------------------%
\subsection{Gabarito:}

\begin{enumerate}

\item
\begin{enumerate}
\item$R(x)=x(1200-12x)$
\item
\begin{center}
  \includegraphics[width=0.5\textwidth]{gabarito_aula1-1b}
\end{center}
\item$x=50$ \\
\end{enumerate}

\item
\begin{enumerate}
\item $3$, se $x \neq 2$
\item $x+4$, se $x \neq 4$ \\
\end{enumerate}

\item
\begin{enumerate}
\item $f(0)=f(2)= \sqrt{3}$ e $f(0,5)=\frac{\sqrt{15}}{2}$
\item $D(f)= \{ x \in \mathbb{R} | -1 \leq x \leq 3\}$
\item Note que o valor mínimo de $f$ é $0$ e o valor máximo é $2$, assim, $Im(f)= [0,2]$
\item $x=-1$ e $x=3$
\item Pelo gráfico podemos notar que $f$ é crescente no intervalo $[-1,1]$ e decrescente no intervalo $[1,3]$
\item $(-1,0)$ e $(3,0)$ são os pontos de mínimo absolutos e $(1,2)$ é o ponto de máximo absoluto
\end{enumerate}



\item
\begin{enumerate}
\item $D(f)=\{ x \in \mathbb{R}| x \neq 0, x \neq -5\}$
\item $D(f)=\{ x \in \mathbb{R}| x \geq 5\}$
\item $D(f)=\{ x \in \mathbb{R}| x \geq \frac{3}{2}, x \neq 4\}$
\item $D(h)=\{ x \in \mathbb{R}| x < 0$ ou $ x > 5\}$   \\
\end{enumerate}

\item
\begin{enumerate}
\item par
\item ímpar
\item par
\item ímpar
\item nem par, nem ímpar
\item nem par, nem ímpar\\
\end{enumerate}

\item
\begin{enumerate}
\item $f$ é nem par, nem ímpar, $g$ é par e $h$ é nem par, nem ímpar
\item $f$ e $h$ são injetoras e $g$ não é injetora
\item
$f$ é sempre crescente, $g$ é decrescente no intevalo $(-\infty,0)$ e crescente no intevalo $(0, \infty)$ e $h$ é sempre decrescente
\begin{figure}[h!]
\centering
\includegraphics[scale=0.55]{gabarito_aula1-6cI}
\caption{Gráfico da função $f(x)=x+2$}
\end{figure}

\begin{figure}[h!]
\centering
\includegraphics[scale=0.55]{gabarito_aula1-6cII}
\caption{Gráfico da função $g(x)=x^2-4$}
\end{figure}

\begin{figure}[h!]
\centering
\includegraphics[scale=0.55]{gabarito_aula1-6cIII}
\caption{Gráfico da função $h(x)=-2x+5$}
\end{figure}
\item $-\frac{4}{3}$
\item $g(h(2))=-3$
\item $ $
  \begin{figure}[h]
  \centering
  \includegraphics[scale=0.55]{gabarito_aula1-6f}
  \caption{Gráfico da função $f(x)=2^{f(x)}$}
  \end{figure}
\end{enumerate}

\item
\begin{enumerate}
\item $D(f)=\{ x \in \mathbb{R}|  x \neq 3\}$
\item $f$ não possui zeros
\item nem par, nem ímpar
\item Intercepta o eixo $y$ em $x=-\frac{1}{3}$ e não intercepta o eixo $x$, visto que não possui zeros
\item $x=-2$ e $x=1$\\
\end{enumerate}

\item
\begin{enumerate}

\item $p=200-l$
\item $A(l)=200l-l^2$
\item $A(75)=9375 m^2$ e $A(150)=7500m^2$
\item $D=(0,200)$, para chegar nesse domínio você deve impor que $A(l)> 0$
\item $ $
        \begin{figure}[h]
        \centering
        \includegraphics[height=6cm , width =10cm]{gabarito_aula1-8e}
        \caption{Gráfico da função $A(l)= 200l-l^2$ }
        \end{figure}
\item É crescente no intervalo $(0,100) $ e decrescente no intervalo $(100,200)$
\item Não, pois a equação $200l-l^2=12000$ não possui solução real
\item $A_{max}=10000m^2$ e as dimensões são $l=p=100$ \\
\end{enumerate}

\item
\begin{enumerate}
    \item $ $
        \begin{figure}[h]
        \centering
        \includegraphics[scale=0.55]{gabarito_aula1-9a}
        \caption{Gráfico da função $g(x)= \sqrt[3]{x+3}$ }
        \end{figure}
    \item $ $
        \begin{figure}[h]
        \centering
        \includegraphics[scale=0.55]{gabarito_aula1-9b}
        \caption{Gráfico da função $h(x)= \sqrt[3]{3x}$ }
        \end{figure}
    \item$ $
        \begin{figure}[h]
        \centering
        \includegraphics[scale=0.55]{gabarito_aula1-9c}
        \caption{Gráfico da função $i(x)= -\sqrt[3]{x}$ }
        \end{figure}
    \item$ $
        \begin{figure}[h]
        \centering
        \includegraphics[scale=0.55]{gabarito_aula1-9d}
        \caption{Gráfico da função $j(x)= 3\sqrt[3]{x}$ }
        \end{figure}
    \item$ $
        \begin{figure}[h]
        \centering
        \includegraphics[scale=0.55]{gabarito_aula1-9e}
        \caption{Gráfico da função $k(x)= 1-\sqrt[3]{x-1}$ }
        \end{figure}
    \item $ $
        \begin{figure}[h]
        \centering
        \includegraphics[scale=0.55]{gabarito_aula1-9f}
        \caption{Gráfico da função $l(x)=\sqrt[3]{2x}-1$ }
        \end{figure}
\end{enumerate}

\item
\begin{enumerate}
\item $f(g(x))=x^2$, $g(f(x))=4x^2$, $f(f(x))=16x$ e $g(g(x))=\frac{x^4}{64}$
\item $f(g(x))=\sqrt{\frac{x}{3}}$, $g(f(x))=\frac{\sqrt{x}}{3}$, $f(f(x))=\sqrt{\sqrt{x}}$ e $g(g(x))=\frac{x}{9}$\\
\end{enumerate}

\item
\begin{enumerate}
\item  $f(g(x))= \frac{1}{x^2-4}$ e $D(f(g))=\{ x \in \mathbb{R}| x \neq -2, x \neq 2 \}$, $g(f(x))=\frac{1}{(x-4)^2}$ e $D(g(f))=\{ x \in \mathbb{R}| x \neq 4\}$ \\
\item  $f(g(-3))= \frac{1}{5}$ e $g(f(7))= \frac{1}{9}$
\end{enumerate}

\item
\begin{enumerate}
    \item $c(d(t)) = 1154,7t^4 - 32091t^3 + 270108t^2 - 648900t + 514422$
    \item O valor interno $d(t)$ calcula o número de casos para um determinado mês ($t$).
    O valor externo $c(d(t))$ calcula o custo total baseado nesse número de casos.
    \item $d(4) \approx783,33$ pessoas. Como o número de pessoas é inteiro devemos calcular o custo para 784 pessoas. Assim $c(784) = \text{R\$ } 482.400,00$
\end{enumerate}
\end{enumerate}

